\documentclass[a4paper,11pt]{article}
\usepackage[margin=0.6in,top=0.7in,bottom=0.6in]{geometry}
\usepackage{fancyhdr}
\usepackage{amsmath}
\usepackage{xcolor}
\usepackage{tikz}
\usepackage{enumitem}
\usepackage[utf8]{inputenc}
\usepackage{tcolorbox}
\usepackage{booktabs}
\usepackage{amssymb}
\usepackage{qrcode}
\usepackage{hyperref}

% Professional Academic Color Scheme
\definecolor{primaryblue}{RGB}{0,51,102}
\definecolor{accentblue}{RGB}{0,102,204}
\definecolor{codebg}{RGB}{248,249,250}
\definecolor{bordercolor}{RGB}{180,180,180}
\definecolor{lightgray}{RGB}{240,240,240}

% Professional Exam Header with QR Code
\newcommand{\examheader}[4]{%
    \noindent
    \begin{tikzpicture}
        \fill[primaryblue] (0,0) rectangle (\textwidth, 1.2cm);
    \node[white] at (0.5\textwidth, 0.8cm) {\Large \textbf{AlMustafa University}};
    \node[white] at (0.5\textwidth, 0.35cm) {\small Cybersecurity Engineering Department};
    \end{tikzpicture}
    
    \vspace{3mm}
    \noindent
    {\large\textbf{Linux Systems Administration - Mid-Term Exam}}
    
    \vspace{2mm}
    \noindent
    \textbf{Instructor:} Dr. Husam Alkinani \hfill \textbf{Exam Sheet:} #2
    
    \noindent
    \textbf{Date:} November 2, 2025 \hfill \textbf{Duration:} #3
    
    \vspace{2mm}
    \noindent\rule{\textwidth}{0.8pt}
    \vspace{2mm}
    
    \noindent
    \textbf{Student Name:} \underline{\hspace{6cm}} \hfill \textbf{Student ID:} \underline{\hspace{3.5cm}}
    
    \vspace{2.5mm}
    \noindent
    \colorbox{lightgray}{\parbox{\dimexpr\textwidth-2\fboxsep\relax}{%
        \small\textbf{Instructions:} Answer all questions clearly. Execute commands in your Linux terminal. Write legibly. \textbf{Total: 20 points}
    }}
    
    % QR Code for Answer Access - Enhanced Size for Camera Scanning
    \vspace{3mm}
    \noindent
    \begin{tcolorbox}[
        colback=white,
        colframe=primaryblue,
        boxrule=1.5pt,
        arc=3pt,
        left=8pt,
        right=8pt,
        top=6pt,
        bottom=6pt,
        width=\textwidth
    ]
    \begin{minipage}{0.65\textwidth}
        \textcolor{primaryblue}{\textbf{Answer Key Access:}} \\
        \small Scan QR code with mobile camera for secure answer verification \\
        \footnotesize \textit{Private GitHub repository access}
    \end{minipage}%
    \begin{minipage}{0.35\textwidth}
    \centering
                    \qrcode[height=3.2cm, padding]{#4}
    \end{minipage}
    \end{tcolorbox}
    
    \vspace{2mm}
}

% Professional Cutting Line
\newcommand{\cuttingline}{%
    \vspace{4mm}
    \noindent
    \begin{tikzpicture}
        \draw[dashed, line width=1pt, gray] (0,0) -- (\textwidth,0);
    \end{tikzpicture}
    \vspace{4mm}
}

% Professional Command Box
\newcommand{\codebox}[1]{%
    \vspace{1mm}
    \noindent
    \begin{tcolorbox}[
        colback=codebg,
        colframe=accentblue,
        boxrule=1pt,
        arc=2pt,
        left=4pt,
        right=4pt,
        top=3pt,
        bottom=3pt,
        width=\textwidth
    ]
    \centering\ttfamily\large\textbf{#1}
    \end{tcolorbox}
    \vspace{1mm}
}

% Professional Answer Box
\newcommand{\answerbox}[1]{%
    \vspace{1mm}
    \noindent
    \begin{tcolorbox}[
        colback=white,
        colframe=bordercolor,
        boxrule=0.8pt,
        arc=1pt,
        height=#1,
        left=6pt,
        right=6pt,
        top=4pt,
        bottom=4pt,
        width=\textwidth
    ]
    \textcolor{gray}{\footnotesize \textit{Write your answer here:}}
    \end{tcolorbox}
}

\pagestyle{fancy}
\fancyhf{}
\fancyfoot[C]{\small \textcolor{gray}{Linux Cybersecurity Lab Exam}}
\renewcommand{\headrulewidth}{0pt}
\renewcommand{\footrulewidth}{0.8pt}

\setlength{\parindent}{0pt}
\setlength{\parskip}{2pt}

\begin{document}

% ================= SHEET A =================
\begin{minipage}[t][\dimexpr0.5\textheight-5mm\relax][t]{\textwidth}
% Sheet A - GitHub Pages Auth URL
\examheader{\_\_\_\_\_\_\_\_\_\_\_\_\_\_\_\_\_\_}{\_\_\_\_\_\_}{A}{https://husam05.github.io/linux-exam-answers/student-auth-77.html?sheet=A}

\textbf{\large\textcolor{primaryblue}{Question 1: Practical Command Execution}} \hfill \textit{(10 points)}

\vspace{1mm}
Execute the following command in your Linux terminal and document the complete output:

\codebox{ls -l /etc/passwd}

\textbf{Command Output:}
\answerbox{26mm}

\textbf{Explanation:} Describe what information this command displays (2-3 lines):

\underline{\hspace{\textwidth}}

\underline{\hspace{\textwidth}}

\vspace{2mm}
\noindent\rule{\textwidth}{0.4pt}
\vspace{2mm}

\textbf{\large\textcolor{primaryblue}{Question 2: Theoretical Knowledge}} \hfill \textit{(10 points)}

\vspace{1mm}
\textbf{Multiple Choice Question:}

In Linux file permissions, what does the notation "rwx" represent?

\begin{enumerate}[label=\textbf{\alph*)}, leftmargin=*, itemsep=2pt, topsep=2pt]
    \item read, write, execute
    \item root, web, extra
    \item run, wait, exit
    \item red, white, expert
\end{enumerate}

\textbf{Your Answer:} \underline{\hspace{2.5cm}}

\textbf{Justification (1 line):} \underline{\hspace{9.5cm}}
\end{minipage}

\cuttingline

% ================= SHEET B =================
\begin{minipage}[t][\dimexpr0.5\textheight-5mm\relax][t]{\textwidth}
\examheader{LINUX LAB EXAM}{B}{30 minutes}{https://husam05.github.io/linux-exam-answers/student-auth-77.html?sheet=B}

\textbf{\large\textcolor{primaryblue}{Question 1: Practical Command Execution}} \hfill \textit{(10 points)}

\vspace{1mm}
Execute the following command in your Linux terminal and document the complete output:

\codebox{whoami}

\textbf{Command Output:}
\answerbox{26mm}

\textbf{Explanation:} Describe the function and purpose of this command (2-3 lines):

\underline{\hspace{\textwidth}}

\underline{\hspace{\textwidth}}

\vspace{2mm}
\noindent\rule{\textwidth}{0.4pt}
\vspace{2mm}

\textbf{\large\textcolor{primaryblue}{Question 2: Theoretical Knowledge}} \hfill \textit{(10 points)}

\vspace{1mm}
\textbf{Multiple Choice Question:}

Which Linux distribution is specifically designed for cybersecurity professionals and penetration testing?

\begin{enumerate}[label=\textbf{\alph*)}, leftmargin=*, itemsep=2pt, topsep=2pt]
    \item Windows Operating System
    \item Kali Linux
    \item Microsoft Office Suite
    \item Google Chrome Browser
\end{enumerate}

\textbf{Your Answer:} \underline{\hspace{2.5cm}}

\textbf{Justification (1 line):} \underline{\hspace{9.5cm}}
\end{minipage}

\newpage

% ================= SHEET C =================
\begin{minipage}[t][\dimexpr0.5\textheight-5mm\relax][t]{\textwidth}
\examheader{LINUX LAB EXAM}{C}{30 minutes}{https://husam05.github.io/linux-exam-answers/student-auth-77.html?sheet=C}

\textbf{\large\textcolor{primaryblue}{Question 1: Practical Command Execution}} \hfill \textit{(10 points)}

\vspace{1mm}
Execute the following command in your Linux terminal and document the complete output:

\codebox{pwd}

\textbf{Command Output:}
\answerbox{26mm}

\textbf{Explanation:} Describe what information this command provides (2-3 lines):

\underline{\hspace{\textwidth}}

\underline{\hspace{\textwidth}}

\vspace{2mm}
\noindent\rule{\textwidth}{0.4pt}
\vspace{2mm}

\textbf{\large\textcolor{primaryblue}{Question 2: Theoretical Knowledge}} \hfill \textit{(10 points)}

\vspace{1mm}
\textbf{Multiple Choice Question:}

What does the Linux command 'sudo' stand for?

\begin{enumerate}[label=\textbf{\alph*)}, leftmargin=*, itemsep=2pt, topsep=2pt]
    \item Super user do
    \item System update download only
    \item Secure user data online
    \item Start up disk operations
\end{enumerate}

\textbf{Your Answer:} \underline{\hspace{2.5cm}}

\textbf{Justification (1 line):} \underline{\hspace{9.5cm}}
\end{minipage}

\cuttingline

% ================= SHEET D =================
\begin{minipage}[t][\dimexpr0.5\textheight-5mm\relax][t]{\textwidth}
\examheader{LINUX LAB EXAM}{D}{30 minutes}{https://husam05.github.io/linux-exam-answers/student-auth-77.html?sheet=D}

\textbf{\large\textcolor{primaryblue}{Question 1: Practical Command Execution}} \hfill \textit{(10 points)}

\vspace{1mm}
Execute the following command in your Linux terminal and document the complete output:

\codebox{df -h}

\textbf{Command Output:}
\answerbox{24mm}

\textbf{Explanation:} Describe what information this command displays (2-3 lines):

\underline{\hspace{\textwidth}}

\underline{\hspace{\textwidth}}

\vspace{2mm}
\noindent\rule{\textwidth}{0.4pt}
\vspace{2mm}

\textbf{\large\textcolor{primaryblue}{Question 2: Theoretical Knowledge}} \hfill \textit{(10 points)}

\vspace{1mm}
\textbf{Multiple Choice Question:}

What does the 'df -h' command display in Linux?

\begin{enumerate}[label=\textbf{\alph*)}, leftmargin=*, itemsep=2pt, topsep=2pt]
    \item Directory file listing
    \item Disk space usage in human-readable format
    \item Date and file information
    \item Default filesystem type
\end{enumerate}

\textbf{Your Answer:} \underline{\hspace{2.5cm}}

\textbf{Justification (1 line):} \underline{\hspace{9.5cm}}
\end{minipage}

\newpage

% ================= SHEET E =================
\begin{minipage}[t][\dimexpr0.5\textheight-5mm\relax][t]{\textwidth}
\examheader{LINUX LAB EXAM}{E}{30 minutes}{https://husam05.github.io/linux-exam-answers/student-auth-77.html?sheet=E}

\textbf{\large\textcolor{primaryblue}{Question 1: Practical Command Execution}} \hfill \textit{(10 points)}

\vspace{1mm}
Execute the following command in your Linux terminal and document the output:

\codebox{uname -a}

\textbf{Command Output:}
\answerbox{23mm}

\textbf{Explanation:} Describe what system information this command provides (2-3 lines):

\underline{\hspace{\textwidth}}

\underline{\hspace{\textwidth}}

\vspace{2mm}
\noindent\rule{\textwidth}{0.4pt}
\vspace{2mm}

\textbf{\large\textcolor{primaryblue}{Question 2: Theoretical Knowledge}} \hfill \textit{(10 points)}

\vspace{1mm}
\textbf{Multiple Choice Question:}

What type of information does the 'uname -a' command display?

\begin{enumerate}[label=\textbf{\alph*)}, leftmargin=*, itemsep=2pt, topsep=2pt]
    \item User account names and passwords
    \item Complete system information including kernel version
    \item Network configuration details
    \item Disk usage and filesystem information
\end{enumerate}

\textbf{Your Answer:} \underline{\hspace{2.5cm}}

\textbf{Justification (1 line):} \underline{\hspace{9.5cm}}
\end{minipage}

\cuttingline

% ================= SHEET F =================
\begin{minipage}[t][\dimexpr0.5\textheight-5mm\relax][t]{\textwidth}
\examheader{LINUX LAB EXAM}{F}{30 minutes}{https://husam05.github.io/linux-exam-answers/student-auth-77.html?sheet=F}

\textbf{\large\textcolor{primaryblue}{Question 1: Practical Command Execution}} \hfill \textit{(10 points)}

\vspace{1mm}
Execute the following command in your Linux terminal and document the output:

\codebox{top}

\textbf{Command Output:}
\answerbox{23mm}

\textbf{Explanation:} Describe what information the 'top' command displays (2-3 lines):

\underline{\hspace{\textwidth}}

\underline{\hspace{\textwidth}}

\vspace{2mm}
\noindent\rule{\textwidth}{0.4pt}
\vspace{2mm}

\textbf{\large\textcolor{primaryblue}{Question 2: Theoretical Knowledge}} \hfill \textit{(10 points)}

\vspace{1mm}
\textbf{Multiple Choice Question:}

What is the primary purpose of the 'top' command in Linux system administration?

\begin{enumerate}[label=\textbf{\alph*)}, leftmargin=*, itemsep=2pt, topsep=2pt]
    \item Display static file system information
    \item Monitor real-time system processes and resource usage
    \item Configure network interfaces
    \item Manage user accounts and permissions
\end{enumerate}

\textbf{Your Answer:} \underline{\hspace{2.5cm}}

\textbf{Justification (1 line):} \underline{\hspace{9.5cm}}
\end{minipage}

\end{document}
